% Generated draft: requires mathematical and source review.
\documentclass[11pt]{article}
\usepackage[margin=25mm]{geometry}
\usepackage{fontspec}
\setmainfont{Latin Modern Roman}
\usepackage{amsmath,amssymb,mathtools,unicode-math}
\setmathfont{Latin Modern Math}
\usepackage{enumitem,needspace,xurl,hyperref}
\setlength{\parindent}{0pt}
\setlength{\parskip}{5pt}
\setlength{\emergencystretch}{3em}
\newcommand{\N}{\mathbb{N}}
\newcommand{\Z}{\mathbb{Z}}
\newcommand{\Q}{\mathbb{Q}}
\newcommand{\R}{\mathbb{R}}
\newcommand{\C}{\mathbb{C}}
\newcommand{\F}{\mathbb{F}}
\newcommand{\A}{\mathbb{A}}
\newcommand{\PP}{\mathbb{P}}
\newcommand{\E}{\mathbb{E}}
\newcommand{\eps}{\varepsilon}
\newcommand{\dd}{\,\mathrm{d}}
\begin{document}
\section*{UnsolvedMath ID: 30002684}
\textbf{Source problem code:} catalogue entry
\subsection{Definitions and mathematical statement}
A Ricci soliton consists of a smooth Riemannian manifold $(M,g)$, a smooth vector field $X$, and a constant $\lambda\in\R$ satisfying
\[
\operatorname{Ric}_g+\frac12\mathcal L_Xg=\lambda g.
\]
Here $\mathcal L_Xg$ is the Lie derivative of the metric. A gradient Ricci soliton additionally satisfies $X=\nabla_g f$ for a smooth function $f$; the supplied question does not explicitly require this. A metric is Einstein if $\operatorname{Ric}_g=cg$ for a constant $c$. The usual non-Einstein interpretation of ``nontrivial'' would exclude such metrics, but the supplied excerpt does not define that term, so this interpretation requires confirmation.

On a complex manifold $(M,J)$, a Hermitian metric satisfies $g(Ju,Jv)=g(u,v)$ for all tangent vectors $u,v$. Its fundamental two-form is $\omega(u,v)=g(Ju,v)$. It is K\"ahler when $\mathrm d\omega=0$, and non-K\"ahler when $\mathrm d\omega\not\equiv0$. Thus, under the non-Einstein interpretation of nontriviality, a candidate Hermitian non-K\"ahler Ricci soliton must satisfy
\[
\operatorname{Ric}_g+\frac12\mathcal L_Xg=\lambda g,
\qquad g(Ju,Jv)=g(u,v),
\qquad \mathrm d\omega\not\equiv0,
\qquad \nexists c\in\R:\ \operatorname{Ric}_g=cg.
\]
These displayed conditions describe candidate objects; they do not resolve the missing global hypotheses of the source question. In particular, the source excerpt does not specify compactness, completeness, boundary conditions, the sign of $\lambda$, whether $X$ must preserve $J$, or a prescribed dimension. Non-K\"ahler Hermitian metrics necessarily have complex dimension at least two.

For reference, a smooth Fano variety is a smooth complex projective variety whose anticanonical line bundle $K_M^{-1}$ is ample. The supplied record does not determine whether its use of ``Fano varieties'' is restricted to this smooth setting or also permits singular varieties, nor does it specify any singularity class or corresponding notion of soliton solution.

\begin{enumerate}
\item Can nontrivial Hermitian non-K\"ahler Ricci solitons be constructed using the techniques referred to in the original discussion, or related techniques successful in the K\"ahler setting?
\item Can methods from partial differential equations establish existence of nontrivial Ricci solitons on a broader class of Fano varieties than the constructions underlying that discussion?
\item Can other special geometries provide constructions of nontrivial Ricci solitons?
\end{enumerate}
The first question concerns existence together with a requested construction method. The second and third are research directions whose comparison classes are not supplied. In particular, the record does not state existence on every Fano variety, does not identify which broader class is intended, and does not explicitly require the solitons in the second or third questions to be non-K\"ahler. A fully quantified existence conjecture cannot be extracted without resolving these qualifications.
\subsection{Short English statement}
The problem asks for constructions of nontrivial Hermitian non-K\"ahler Ricci solitons, broader existence results on Fano varieties using partial differential equations, and constructions arising from other special geometries. The source leaves the global hypotheses, the meaning of nontriviality, and the intended comparison classes unspecified. This is a statement-only draft requiring human mathematical review.
\subsection{Sources}
\begin{itemize}
\item Local UnsolvedMath registry: supplied record with permanent ID \texttt{30002684}, exact source problem code \texttt{catalogue entry}, title ``Non-K\"ahler Ricci Solitons and Broader Existence Methods,'' in the Oberwolfach Reports Open Problems collection. The proposed year is recorded as 2014; no individual proposer is supplied.
\item Supplied source citation: \emph{Mini-Workshop: Einstein Metrics, Ricci Solitons and Ricci Flow under Symmetry Assumptions (2015)}, Oberwolfach Reports, pp.~2529--2568, DOI \texttt{10.4171/owr/2014/45}. Supplied source URL and literature-source link: \url{https://doi.org/10.4171/owr/2014/45}. The literature-source label is ``Oberwolfach Report 45/2014.'' The supplied citation's 2015 designation differs from the report's 2014 identifier and the recorded proposed year; this discrepancy has not been resolved.
\item Additional supplied literature source: \emph{Equivariant Solutions to the Prescribed Q-Curvature Equation}, \url{https://doi.org/10.1007/s12220-024-01554-4}. Its relevance to the precise Ricci soliton questions is not explained in the record.
\item Supplied literature assessment dated exactly \texttt{2026-08-21}: the record labels the problem open and describes the targets as existence beyond known K\"ahler and symmetric constructions. It reports that the literature still describes compact non-Einstein Ricci soliton examples as K\"ahler and treats the non-K\"ahler and broader Fano/PDE questions as unresolved. This is supplied dated curation, not an independently verified assessment. The background expressly requests independent verification of time-sensitive status and cited claims.
\item The supplied statement status is \texttt{corrected\_verified}, accompanied by the claim that the clean statement was checked against the official Oberwolfach report and preserves its hypotheses, notation, and open target. This verification claim is part of the supplied metadata; no source consultation or independent verification is asserted here.
\end{itemize}
\end{document}
